% !TEX root = ../matan.tex

\section{Скалярное произведение. Примеры. Неравенство Коши--Буняковского--Шварца. Норма в пространстве со скалярным произведением}

\subsection*{Скалярное произведение}

\begin{Definition}{Скалярное произведение}{}
	Пусть $V$ --- линейное пространство над $\RR$. \textbf{Скалярным
	произведением} называется отображение
	$\langle\cdot,\cdot\rangle:V\times V\to\RR$, удовлетворяющее для всех
	$x,y,z\in V$ и $\lambda\in\RR$ свойствам:
	\begin{enumerate}
		\item \emph{положительная определённость}:
		$\langle x,x\rangle\ge0$, причём $\langle x,x\rangle=0\iff x=0$;
		\item \emph{симметрия}: $\langle x,y\rangle=\langle y,x\rangle$;
		\item \emph{линейность по первому аргументу}:
		$\langle\lambda x,y\rangle=\lambda\langle x,y\rangle$ и
		$\langle x+y,z\rangle=\langle x,z\rangle+\langle y,z\rangle$.
	\end{enumerate}
	(Из симметрии следует линейность и по второму аргументу.) Пространство со
	скалярным произведением называют \textbf{евклидовым} (предгильбертовым); если
	оно полно относительно порождённой нормы --- \textbf{гильбертовым}.
\end{Definition}

\begin{Example}{Примеры скалярных произведений}{}
	\begin{enumerate}
		\item В $\RR^n$: $\langle x,y\rangle=\sum_{k=1}^{n} x_k y_k$ ---
		стандартное скалярное произведение.
		\item В $\ell^2$: $\langle x,y\rangle=\sum_{k=1}^{\infty} x_k y_k$ (ряд
		сходится абсолютно по неравенству Гёльдера с $p=q=2$).
		\item В $C[a,b]$: $\langle f,g\rangle=\int_a^b f(x)g(x)\,dx$. Все аксиомы
		выполнены; положительная определённость использует непрерывность: если
		$\int_a^b f^2=0$ и $f$ непрерывна, то $f\equiv0$.
	\end{enumerate}
\end{Example}

\subsection*{Норма, порождённая скалярным произведением}

В пространстве со скалярным произведением полагают
\[
\|x\|:=\sqrt{\langle x,x\rangle}\,;
\]
это определение корректно, так как $\langle x,x\rangle\ge0$. Докажем, что
$\|\cdot\|$ --- норма; единственное нетривиальное свойство (неравенство
треугольника) выводится из неравенства Коши--Буняковского--Шварца.

\begin{Theorem}{Неравенство Коши--Буняковского--Шварца}{}
	Для любых $x,y$ в пространстве со скалярным произведением
	\[
	|\langle x,y\rangle|\le\|x\|\,\|y\|,
	\]
	причём равенство достигается тогда и только тогда, когда $x,y$ линейно
	зависимы.
\end{Theorem}

\begin{proof}
	Если $y=0$, то обе части равны нулю (так как
	$\langle x,0\rangle=\langle x,0\cdot 0\rangle=0\cdot\langle x,0\rangle=0$ и
	$\|0\|=0$), и неравенство верно. Пусть теперь $y\ne0$, тогда
	$\|y\|^2=\langle y,y\rangle>0$. Для любого $\lambda\in\RR$, пользуясь
	билинейностью и симметрией,
	\[
	0\le\langle x+\lambda y,\,x+\lambda y\rangle
	=\|x\|^2+2\lambda\langle x,y\rangle+\lambda^2\|y\|^2.
	\]
	Это квадратный трёхчлен относительно $\lambda$ с положительным старшим
	коэффициентом $\|y\|^2>0$, неотрицательный при всех $\lambda$. Значит, его
	дискриминант неположителен:
	\[
	\bigl(2\langle x,y\rangle\bigr)^2-4\|y\|^2\|x\|^2\le0
	\quad\Longrightarrow\quad
	\langle x,y\rangle^2\le\|x\|^2\|y\|^2.
	\]
	Извлекая корень, получаем $|\langle x,y\rangle|\le\|x\|\,\|y\|$.

	\emph{Случай равенства.} Равенство в КБШ при $y\ne0$ означает, что
	дискриминант равен нулю, то есть трёхчлен имеет (двойной) корень
	$\lambda_0$, при котором $\langle x+\lambda_0 y,x+\lambda_0 y\rangle=0$. По
	положительной определённости это равносильно $x+\lambda_0 y=0$, то есть
	$x=-\lambda_0 y$ --- векторы линейно зависимы. Обратно, при $x=\mu y$
	обе части равны $|\mu|\,\|y\|^2$.
\end{proof}

\begin{Theorem}{Норма из скалярного произведения}{}
	Функция $\|x\|=\sqrt{\langle x,x\rangle}$ является нормой на $V$.
\end{Theorem}

\begin{proof}
	\emph{Невырожденность}: $\|x\|=\sqrt{\langle x,x\rangle}\ge0$, и $\|x\|=0
	\iff\langle x,x\rangle=0\iff x=0$ по положительной определённости.
	\emph{Однородность}:
	$\|\lambda x\|^2=\langle\lambda x,\lambda x\rangle=\lambda^2\langle x,x\rangle
	=\lambda^2\|x\|^2$, откуда $\|\lambda x\|=|\lambda|\,\|x\|$.
	\emph{Неравенство треугольника}: пользуясь КБШ,
	\[
	\|x+y\|^2=\langle x+y,x+y\rangle
	=\|x\|^2+2\langle x,y\rangle+\|y\|^2
	\le\|x\|^2+2\|x\|\,\|y\|+\|y\|^2
	=\bigl(\|x\|+\|y\|\bigr)^2.
	\]
	Так как обе части неотрицательны, $\|x+y\|\le\|x\|+\|y\|$.
\end{proof}

\begin{Remark}{Тождество параллелограмма}{}
	Норма, порождённая скалярным произведением, удовлетворяет
	\textbf{тождеству параллелограмма}
	\[
	\|x+y\|^2+\|x-y\|^2=2\|x\|^2+2\|y\|^2,
	\]
	которое проверяется раскрытием скобок через $\langle\cdot,\cdot\rangle$.
	Это тождество выделяет «евклидовы» нормы среди всех норм: например, нормы
	$\|\cdot\|_p$ при $p\ne2$ ему не удовлетворяют и потому не порождаются
	никаким скалярным произведением. Само скалярное произведение
	восстанавливается по норме поляризационным тождеством
	$\langle x,y\rangle=\tfrac14\bigl(\|x+y\|^2-\|x-y\|^2\bigr)$.
\end{Remark}

\begin{Example}{КБШ в конкретных пространствах}{}
	В $\RR^n$ неравенство КБШ принимает вид
	$\bigl(\sum_{k} x_k y_k\bigr)^2\le\bigl(\sum_k x_k^2\bigr)\bigl(\sum_k y_k^2\bigr)$,
	а в $C[a,b]$ ---
	$\bigl(\int_a^b fg\,dx\bigr)^2\le\bigl(\int_a^b f^2\,dx\bigr)\bigl(\int_a^b g^2\,dx\bigr)$
	(интегральное неравенство Коши--Буняковского).
\end{Example}
